% Четыре четверти коврика и сдвиг номеров клеток
\begin{tikzpicture}[x=1.0cm, y=1.0cm, font=\sffamily]
  \foreach \q/\x/\y/\s in {1/0/0/0, 2/6/0/6, 3/0/6/72, 4/6/6/78}{
    \fill[accentsoft] (\x,\y) rectangle (\x+6,\y+6);
    \draw[accent, very thick] (\x,\y) rectangle (\x+6,\y+6);
    \node[text=accent, font=\large\bfseries] at (\x+3,\y+4.4) {Четверть \q};
    \node[text=ink, font=\small] at (\x+3,\y+3.4) {левая нижняя клетка: \s};
    \node[text=inkgray, font=\small] at (\x+3,\y+2.4) {сдвиг номеров: \ifnum\s=0 нет\else $+\s$\fi};
  }
  \node[text=inkgray, font=\scriptsize] at (6,-0.4) {клетка 0 внизу слева, номера растут вправо и вверх};
\end{tikzpicture}
